\documentclass{article}
\usepackage{amsmath,amssymb,amsthm,mathtools,mathrsfs}

\begin{document}
\maketitle
\section{Independent completed-phase quotient review}\label{independent-completed-phase-quotient-review}

Date: 2026-09-13. Scope: the polynomial-density, completed-relation, kernel, and direct-integral assertions in H64. No source file was edited; no scans, program tests, Lean runs, or commutator calculations were performed.

\subsection{Reviewed text and conclusion}\label{reviewed-text-and-conclusion}

The review uses the following complete supplied sources. The original relative source names below are retained for provenance; their delivered locations and full source hashes are recorded in the companion source manifest.

The relevant fixed text is:

\begingroup\raggedright
\begin{itemize}
\tightlist
\item
  \nolinkurl{web_holonomy_descent_delivery/Tau_Holonomy_Descent_Control/NOTE.md}, lines 67--90: the complete selected orders, original cyclic polynomial, and Taylor unit; lines 166--177: H17 and positivity for tensor degree at least two; lines 703--715: H64, the asserted kernel, and the almost-everywhere direct-integral statement.
\item
  \nolinkurl{web_holonomy_descent_delivery/Tau_Holonomy_Descent_Control/SOURCE_REVIEW.md}: the source-reading scope was read. This review adds no primary-source scan claim.
\item
  \nolinkurl{web_periodized_source_delivery/Tau_Periodized_Source_Control/RESEARCH_NOTE.md}, lines 88--113: the full-order packet and original cyclic source; lines 564--584: P53--P54 and the exponential-tail density argument; lines 586--612: P55--P57, the target metric, and the closure calculation; line 633: the positive-support qualification for tensor degree one.
\item
  \nolinkurl{web_cyclic_sum_delivery/Tau_Cyclic_Sum_Control/RESEARCH_NOTE.md}, lines 31--34: the full-order packet and unit \(j_h(g/h)\); lines 54--70: the cyclic minimal polynomial and inclusion.
\item
  \nolinkurl{web_arithmetic_endpoint_delivery/Tau_Arithmetic_Endpoint_Bounds/RESEARCH_NOTE.md}, lines 95--105: the full-order packet and the entire quotient defining the actual density.
\end{itemize}
\endgroup

\textbf{Conclusion.} H64 and its kernel are correct for the stated original full-order arithmetic source at every integer tensor degree \(k \ge 1\). For a general nonnegative sampled density, its right side must include only positive-mass sampled roots. At \(k = 1\), this qualification makes no change to H64 for the actual source: the ideals \((\chi_{h,1})=(h)\) agree, and every sampled root has positive mass by full-order cancellation with the original leading coefficient \(a_h\) retained. The direct-integral target is zero, with source kernel all of \(E\), in both the general positive-support formulation and its arithmetic specialization. The proof below supplies the explicit metric and the density argument that are abbreviated in H64.

\subsection{1. Coordinates, mass, and exponential moments}\label{coordinates-mass-and-exponential-moments}

Fix the original integer \(k \ge 1\), period \(L > 0\), phase \(0 \le \theta < 2 \pi\), and the original nonempty full-order packet \(h\). Throughout,

\[
c=k/2,\qquad
u_{n,\theta}=\frac{2\pi n+\theta}{L},\qquad
S_{n,\theta}=c+i u_{n,\theta},\qquad n\in\mathbb Z,
\tag{CQ1}
\]

\[
v_h=g/h,\qquad g=2\xi,\qquad
w_h(t)=\frac{|v_h(1/2+it)|^2}{2\pi},\qquad
m_{h,k}=w_h^{*k},\qquad
\nu_{n,\theta}=\frac{2\pi}{L}m_{h,k}(u_{n,\theta}).
\tag{CQ2}
\]

No mass is divided out. Put

\[
I_\theta=\{n\in\mathbb Z:\nu_{n,\theta}>0\},\qquad
\mathcal H_\theta=\ell^2(I_\theta,\nu_\theta).
\tag{CQ3}
\]

The notation \(\ell^2(\nu)\) on all integers means this same Hilbert space after quotienting the zero seminorm coordinates. A zero-mass coordinate is not an evaluation functional on that Hilbert space.

The predecessor's actual-source input, explicitly recorded immediately before P54, is that for each \(0 < b < \pi/2\),

\[
f_b(t)=e^{b|t|}w_h(t)\in L^1(\mathbb R)\cap L^\infty(\mathbb R).
\tag{CQ4}
\]

This is an inherited estimate for the specified entire quotient, not an extra assumption about a new density. Its consequence for sampled measures is proved here. In the convolution integral, \(|t_1+\cdots+t_k| \le |t_1|+\cdots+|t_k|\), so

\[
e^{b|u|}m_{h,k}(u)
\le f_b^{*k}(u)
\le C_{b,h,k},\qquad
C_{b,h,k}:=\|f_b\|_\infty\|f_b\|_1^{k-1}.
\tag{CQ5}
\]

The second inequality follows by integrating one bounded factor and the remaining \(k-1\) factors in \(L^1\). It includes \(k=1\). For any \(0 \le a < b\), the literal sum satisfies

\[
\begin{aligned}
\sum_{n\in I_\theta}\nu_{n,\theta}e^{a|u_{n,\theta}|}
&\le\frac{2\pi C_{b,h,k}}L
\sum_{n\in\mathbb Z}e^{-(b-a)|2\pi n+\theta|/L}\\
&=\frac{2\pi C_{b,h,k}}L
\frac{e^{-(b-a)\theta/L}+e^{-(b-a)(2\pi-\theta)/L}}
{1-e^{-2\pi(b-a)/L}}<\infty.
\end{aligned}
\tag{CQ6}
\]

The two series use \(n \ge 0\) and \(n \le -1\), respectively. Formula CQ6 also includes \(\theta=0\) without deleting its zeroth term.

Let \(\chi(S)=\sum_{j=0}^q \chi_j S^j\) be the literal monic original polynomial, of degree \(q>0\). For real \(u\),

\[
|\chi(c+iu)|\le
C_\chi(1+|u|)^q,\qquad
C_\chi=\sum_{j=0}^q|\chi_j|\max(1,c)^j.
\tag{CQ7}
\]

Choose \(0 < a_1 < a_2 < b\). The continuous function \((1+t)^{2q} \exp(-(a_2-a_1)t)\) is bounded on \(t \ge 0\), since it tends to zero at infinity. Multiplying CQ7 by the exponential and using CQ6 proves

\[
\sum_{n\in I_\theta}
|\chi(S_{n,\theta})|^2\nu_{n,\theta}
e^{a_1|u_{n,\theta}|}<\infty.
\tag{CQ8}
\]

Thus the original discrete measure and its \(|\chi|^2\) multiple each have a finite positive exponential moment. All their polynomial moments exist.

\subsection{2. A complete polynomial-density proof on the shifted lattice}\label{a-complete-polynomial-density-proof-on-the-shifted-lattice}

The following argument applies to either measure. Retain the shifted lattice CQ1, a subset \(I\) of its indices, and positive masses \(\mu_n\) such that

\[
\sum_{n\in I}\mu_ne^{a|u_{n,\theta}|}<\infty
\quad\text{for some }a>0.
\tag{CQ9}
\]

Suppose \(f\) in \(\ell^2(I,\mu)\) is orthogonal to all polynomial evaluation vectors. Since \(u=(S-c)/i\), the evaluations of polynomials in \(S\) and in the real coordinate \(u\) span the same subspace. Conjugating the orthogonality equations if necessary gives

\[
\sum_{n\in I}f_n\mu_nu_{n,\theta}^{j}=0
\quad(j=0,1,2,\ldots).
\tag{CQ10}
\]

Define

\[
F(z)=\sum_{n\in I}f_n\mu_ne^{zu_{n,\theta}},
\qquad |\Re z|<a/2.
\tag{CQ11}
\]

For \(0 \le t < a/2\), Cauchy--Schwarz gives

\[
\sum_{n\in I}|f_n|\mu_ne^{t|u_{n,\theta}|}
\le\|f\|_{\ell^2(\mu)}
\left(\sum_{n\in I}\mu_ne^{2t|u_{n,\theta}|}\right)^{1/2}<\infty.
\tag{CQ12}
\]

This proves locally uniform absolute convergence on the open strip in CQ11. For every derivative order \(j\) and every smaller closed strip, choose \(t'\) strictly between its width and \(a/2\). The bound

\(\sup_{x \ge 0} x^j \exp(-(t'-t)x) < \infty\)

and CQ12 at \(t'\) prove locally uniform absolute convergence of the derivative series. Hence \(F\) is holomorphic and CQ10 says that all its derivatives at zero vanish. Its convergent Taylor series is zero on a disk about zero. Holomorphic uniqueness on the connected strip therefore gives \(F(z)=0\) on that entire strip.

There is no need to appeal to an unproved general Fourier-measure uniqueness statement here. CQ12 at \(t=0\) allows termwise integration, and for each \(m \in I\),

\[
\begin{aligned}
0&=\frac1L\int_0^L F(it)e^{-it u_{m,\theta}}\,dt\\
&=\sum_{n\in I}f_n\mu_n
\frac1L\int_0^L e^{2\pi i(n-m)t/L}\,dt
=f_m\mu_m.
\end{aligned}
\tag{CQ13}
\]

Because \(\mu_m>0\), every coordinate \(f_m\) vanishes. The orthogonal complement of the polynomial span is zero, so that span is dense. CQ6 proves density in \(\mathcal H_\theta\); CQ8 proves density in the Hilbert space for the positive atoms of \(|\chi|^2 \nu_\theta\).

\subsection{3. The closed relation space, exact map, metric, and kernel}\label{the-closed-relation-space-exact-map-metric-and-kernel}

Define

\[
Z_\theta=\{n\in I_\theta:\chi(S_{n,\theta})=0\},\qquad
p_\theta(S)=\prod_{n\in Z_\theta}(S-S_{n,\theta}),
\tag{CQ14}
\]

where an empty product is \(1\). The sample points are distinct, so \(Z_\theta\) has at most \(q\) members and \(p_\theta\) is a squarefree divisor of \(\chi\).

Let \(\mathcal R_\theta\) be the closure in \(\mathcal H_\theta\) of the evaluation vectors of \(\chi \mathbb C[S]\). Every such vector vanishes at all indices in \(Z_\theta\). Each of those coordinates is continuous, since

\[
|f_n|\le\nu_{n,\theta}^{-1/2}\|f\|_{\mathcal H_\theta}
\qquad(n\in Z_\theta).
\tag{CQ15}
\]

It follows that \(\mathcal R_\theta\) is contained in the subspace of vectors vanishing on \(Z_\theta\). For the reverse inclusion, let \(f\) vanish there and set

\[
g_n=\frac{f_n}{\chi(S_{n,\theta})}
\quad(n\in I_\theta\setminus Z_\theta).
\tag{CQ16}
\]

This vector is square-integrable for the weighted masses

\(\mu_n = |\chi(S_{n,\theta})|^2 \nu_{n,\theta}\),

because its squared norm is exactly \(\sum_{I_\theta\setminus Z_\theta} |f_n|^2 \nu_{n,\theta}\). By Section 2 there are polynomials \(P_j\) tending to \(g\) in that weighted norm. Multiplication by the literal \(\chi\) gives

\[
\|\chi P_j-f\|_{\mathcal H_\theta}^2
=\sum_{n\in I_\theta\setminus Z_\theta}
|P_j(S_{n,\theta})-g_n|^2
|\chi(S_{n,\theta})|^2\nu_{n,\theta}\longrightarrow0.
\tag{CQ17}
\]

At indices in \(Z_\theta\), both terms on the left vanish. Consequently

\[
\mathcal R_\theta
=\{f\in\mathcal H_\theta:f_n=0\text{ for all }n\in Z_\theta\}.
\tag{CQ18}
\]

Restriction to \(Z_\theta\) therefore induces the exact isometric isomorphism

\[
\mathcal H_\theta/\mathcal R_\theta
\xrightarrow{\sim}
\left(\mathbb C^{Z_\theta},
\ \|(v_n)\|^2=\sum_{n\in Z_\theta}\nu_{n,\theta}|v_n|^2\right).
\tag{CQ19}
\]

The quotient norm is exactly the displayed norm: the coordinate projection onto \(Z_\theta\) is orthogonal, and the complementary projection has range \(\mathcal R_\theta\).

Since adding \(\chi Q\) to a polynomial changes its evaluation vector by an element of \(\mathcal R_\theta\), there is a well-defined linear comparison

\[
T_\theta:E=\mathbb C[S]/(\chi)
\longrightarrow\mathcal H_\theta/\mathcal R_\theta,
\qquad
[P]\longmapsto(P(S_{n,\theta}))_{n\in Z_\theta}.
\tag{CQ20}
\]

It is surjective. When \(Z_\theta\) is nonempty, the literal Lagrange polynomials

\[
L_n(S)=\prod_{m\in Z_\theta\setminus\{n\}}
\frac{S-S_{m,\theta}}{S_{n,\theta}-S_{m,\theta}}
\tag{CQ21}
\]

map to its coordinate basis; when it is empty, the target is zero. A polynomial vanishes at every distinct point indexed by \(Z_\theta\) if and only if it is divisible by their squarefree product. Hence the exact kernel is

\[
\ker T_\theta=(p_\theta)/(\chi)\subseteq E.
\tag{CQ22}
\]

This retains the multiplicities in \(\chi\). In the primary decomposition of \(E\), each unsampled primary block lies in the kernel; at a sampled root of multiplicity \(m\), the kernel is the maximal ideal \((S-\rho)/(S-\rho)^m\), while its scalar value survives. CQ20 is a linear map to a Hilbert quotient; no multiplication of arbitrary square-summable sequences is required.

\subsection{4. Why the original arithmetic source needs no extra positivity hypothesis}\label{why-the-original-arithmetic-source-needs-no-extra-positivity-hypothesis}

For \(k \ge 2\), the actual entire nonzero function \(v_h\) has isolated zeros. Thus \(w_h(t)>0\) off a discrete subset of the real axis. For fixed \(u\), both \(w_h(t)\) and \(w_h(u-t)\) are positive for almost every \(t\); their product is integrable by boundedness and integrability from CQ4. Its integral is strictly positive. Thus \(m_{h,2}(u)>0\) for every real \(u\). Inductively, convolving the everywhere-positive \(m_{h,k-1}\) with the almost-everywhere-positive nonzero \(w_h\) gives \(m_{h,k}(u)>0\). CQ5 makes the integrals finite. Therefore \(I_\theta=\mathbb Z\) for \(k \ge 2\) and CQ14 is exactly the index set displayed in H64.

For \(k=1\), the original cyclic algebra is \(\mathbb C[S]/(h)\), with generator multiplication by \(S\). Its annihilator ideal is exactly \((h)\): \(h(M_S)=0\), and any polynomial annihilating \(M_S\) also annihilates \(1\), so is divisible by the unchanged \(h\). Consequently \((\chi_{h,1})=(h)\). If \(\chi_{h,1}\) denotes the monic minimal polynomial and \(a_h\) is the original leading coefficient of \(h\), then \(\chi_{h,1}=h/a_h\); this identifies the ideal generator and does not replace the original source \(v_h=g/h\). Equality \(\chi_{h,1}=h\) is asserted only when the source separately fixes \(a_h=1\).

Write the unchanged packet as

\[
h(S)=a_h\prod_{\rho\in Z_h}(S-\rho)^{m_\rho},
\qquad a_h\ne0,\qquad m_\rho=\operatorname{ord}_\rho g.
\tag{CQ23}
\]

For a selected root \(\rho\), full-order removal gives the exact nonzero value

\[
v_h(\rho)=
\frac{g^{(m_\rho)}(\rho)}{h^{(m_\rho)}(\rho)}=
\frac{g^{(m_\rho)}(\rho)}
{a_hm_\rho!\prod_{\sigma\in Z_h\setminus\{\rho\}}
(\rho-\sigma)^{m_\sigma}}\ne0.
\tag{CQ24}
\]

This follows by factoring \((S-\rho)^{m_\rho}\) from both \(g\) and the unchanged \(h\); the numerator derivative, original leading coefficient \(a_h\), and every denominator factor are nonzero by their definitions. If a sample satisfies \(S_{n,\theta}=\rho\), then \(\Re \rho=1/2\) and

\[
\nu_{n,\theta}
=\frac{2\pi}L w_h(\Im\rho)
=\frac1L|v_h(\rho)|^2
=\frac{|g^{(m_\rho)}(\rho)|^2}
{L|a_h|^2(m_\rho!)^2\prod_{\sigma\in Z_h\setminus\{\rho\}}|\rho-\sigma|^{2m_\sigma}}>0.
\tag{CQ25}
\]

Thus every root appearing in H64 has positive mass also at \(k=1\). Other zero-mass lattice coordinates are discarded by the Hilbert-space definition CQ3, but none removes a summand from H64. The full-order hypothesis at NOTE.md line 67 and the preserved unit at lines 85--90 already supply the needed fact. A useful exposition repair is to define \(I_\theta\) and then state CQ25; it does not change H64's conclusion for this source.

For comparison, outside this arithmetic domain the omission of positive support can change the theorem. With \(\nu_0=0\), \(\nu_n=e^{-|n|}\) for \(n \ne 0\), and \(\chi(S)=S-S_{0,\theta}\), Section 3 gives zero completed quotient, whereas a sum over all sampled polynomial roots would give one copy of \(\mathbb C\). This example establishes why the general formula must use CQ14; it is not a counterexample to the full-order arithmetic result.

\subsection{5. The direct-integral target and its exact zero map}\label{the-direct-integral-target-and-its-exact-zero-map}

Retain the fixed polynomial \(\chi\) and \(c=k/2\). Let

\[
\mathcal F=
\{\theta\in[0,2\pi):
\theta\equiv L\Im\rho\pmod{2\pi}
\text{ for some root }\rho\text{ of }\chi
\text{ with }\Re\rho=c\}.
\tag{CQ26}
\]

Each such root specifies exactly one phase in the chosen half-open interval, and then exactly one integer \(n\) with \(\rho=S_{n,\theta}\). Therefore \(\mathcal F\) is finite, with at most \(q\) elements. For \(\theta\) outside \(\mathcal F\), the index set \(Z_\theta\) is empty, so CQ19 is the zero Hilbert space.

A concrete measurable realization prevents ambiguity about the field of targets. Map CQ19 into the fixed Hilbert space \(\ell^2(\mathbb Z)\) by

\[
(v_n)_{n\in Z_\theta}\longmapsto
\left(\mathbf1_{n\in Z_\theta}\sqrt{\nu_{n,\theta}}v_n\right)_{n\in\mathbb Z}.
\tag{CQ27}
\]

This is an isometry onto the range of the diagonal orthogonal projection with diagonal entries \(\mathbf1_{\nu_{n,\theta}>0} \mathbf1_{\chi(S_{n,\theta})=0}\). Those entries are measurable: the first condition is measurable for the original continuous density, and the second is the zero set of a continuous function of \(\theta\). Hence these ranges define a measurable Hilbert field. Their direct integral has norm

\[
\int_0^{2\pi}\sum_{n\in Z_\theta}
\nu_{n,\theta}|v_n(\theta)|^2\,\frac{d\theta}{2\pi}.
\tag{CQ28}
\]

Every section is zero outside the finite set \(\mathcal F\); its norm in CQ28 is zero. The direct integral is therefore the zero Hilbert space. For each \([P] \in E\), the evaluation sections of CQ20 are measurable through CQ27 and supported in \(\mathcal F\). They consequently define the zero class in the direct integral. The induced map and its exact kernel are

\[
E\longrightarrow
\int_{[0,2\pi)}^\oplus
(\mathcal H_\theta/\mathcal R_\theta)\,\frac{d\theta}{2\pi}
=0,
\qquad \ker=E.
\tag{CQ29}
\]

Individual exceptional fibres can remain nonzero and carry the positive masses in CQ19; a finite set has zero measure for the actual phase measure \(d\theta/(2 \pi)\). CQ29 proves precisely the completed-target claim at NOTE.md line 715. It does not change the source isometry H13 or the finite quotient H24.

On a retained split support label \(\ell\), the actual induced map is \((\ell,v) \longmapsto (\ell,0)\), while external absence maps to external absence. This is the lift of CQ29 by the stated H4; its coefficient kernel is all of \(E\), exactly as CQ29 records.

\subsection{Review disposition}\label{review-disposition}

Accept H64, its displayed squarefree kernel, and the zero direct-integral conclusion for the original full-order packet. The recommended addition is explanatory: define the positive support and retained metric by CQ3, CQ14, CQ19; explicitly record the degree-one identity and positive sampled-root value CQ25. The unrestricted all-atom version is not justified for arbitrary other nonnegative densities. No source edit or mathematical test result is claimed by this review.

\end{document}
